\section{Introduction}
\label{sec:introduction}

When a language model generates the statement ``The Earth orbits the Sun,'' is it drawing on the same internal machinery as when it generates ``Democracy is the best form of government''? The first is an epistemic belief---a factual claim grounded in evidence---while the second is a non-epistemic belief---an opinion reflecting values and preferences. Whether large language models (\llms) internally differentiate between these two categories of belief has direct consequences for AI safety: if factual claims and opinions occupy distinct representational subspaces, targeted interventions could selectively improve factual reliability without distorting the model's handling of subjective content.

Recent work in mechanistic interpretability has demonstrated that \llms encode surprisingly rich semantic structure in their internal activations. \citet{marks2023geometry} showed that truth value is linearly decodable from middle layers of LLaMA models, while \citet{zhu2024language} found that belief states of self and others are represented in distinct attention heads. \citet{bortoletto2024brittle} extended this to show that belief representations are low-dimensional and can be steered through activation addition. However, all of these studies treat ``belief'' or ``truth'' as a monolithic category. None has asked whether the model's internal representations reflect the fundamental distinction---well-established in philosophy and psychology \citep{vesga2025psychological}---between epistemic beliefs (knowledge claims with evidential basis) and non-epistemic beliefs (opinions, preferences, and values).

We address this gap with a systematic layer-wise probing analysis across the GPT-2 model family (117M to 1.5B parameters). We extract residual stream activations at every layer for three categories of statements---epistemic true/false claims, non-epistemic commonly-held/uncommon opinions, and mixed epistemic/non-epistemic statements---and apply mass-mean probing \citep{marks2023geometry} with 5-fold cross-validation. We find three key results:

\begin{itemize}[leftmargin=*,itemsep=0pt,topsep=0pt]
    \item All four GPT-2 models achieve \textbf{99.2--100\% accuracy} at classifying epistemic versus non-epistemic statements from middle layers onward, with selectivity scores of 0.45--0.49 confirming genuine encoding rather than probe memorization.
    \item Non-epistemic agreement is \textbf{more linearly decodable} (87--95\% peak accuracy) than epistemic truth across mixed domains (54--59\%), suggesting \llms encode sentiment and normativity more strongly than factual truth value.
    \item Epistemic truth probing accuracy \textbf{scales with model size} (53.5\% for 117M to 59.0\% for 1.5B on mixed domains; 69.5\% to 85.5\% for geographic facts), while the layer-wise profile of belief encoding shifts systematically with depth.
\end{itemize}

These findings have practical implications for hallucination detection. If epistemic and non-epistemic statements occupy distinct activation subspaces, a detector could first classify a model's output by belief type and then apply type-specific truth verification. Our results also contribute to the ongoing debate about whether \llms develop genuine semantic understanding or merely capture surface-level statistical patterns \citep{marks2023geometry,liu2023cognitive}: the near-perfect type classification accuracy, combined with the modest truth probing accuracy, suggests that \llms encode \emph{how different belief types sound} more robustly than \emph{what is actually true}.

The remainder of this paper is organized as follows. \Secref{sec:related_work} reviews related work on probing and belief representations. \Secref{sec:methodology} describes our experimental methodology. \Secref{sec:results} presents the main results. \Secref{sec:discussion} discusses implications and limitations, and \secref{sec:conclusion} concludes.
